%% kapa: exam=1 round=22 year=2011 period=1 form=A subject=civil_law subject_ko=민법 seq=24 number=24
%% source: Q-Net 감정평가사 기출문제 (cst00309) artlSeq=5036790 file=extracted/1008/5036790/2011년 제22회 감정평가사 1차시험 기출문제/1교시.hwp
%% conversion: kordoc 4.24.51 -> tools/gen1.py (auto); review=proofread-sonnet; flags=none
\begin{question}{24}
대리에 관한 설명으로 옳지 않은 것은? (다툼이 있으면 판례에 의함)
\begin{choices}
  \item 대리인이 파산선고를 받으면 대리권이 소멸한다.
  \item 대리인에게 의사표시의 하자가 있는 때에는 그 효과로서 생기는 취소권은 대리인이 아니라 본인에게 귀속한다.
  \item 부동산 입찰절차에서 동일한 물건에 관하여 한 사람이 동시에 다른 두 사람의 대리인으로서 한 입찰행위는 무효이다.
  \item 대리인이 수인인 때에는 특별한 사정이 없는 한 공동으로 대리하여야 한다.
  \item 본인이 대리인에게 채무를 진 경우, 대리인은 본인의 허락이 없더라도 그 채무를 자신에게 변제할 수 있다.
\end{choices}
\end{question}
\answer{4}
